\rhead{CY401: Ethical Hacking \& Penetration Testing}
\phantomsection 
\section*{Ethical Hacking \& Penetration Testing (CY401)}
\label{major:CS_4}

\noindent \small{\hyperref[main:scheme]{$\leftarrow$ Back to Scheme}} \vspace{0.5cm}

\noindent \textbf{Credits:} 3 (2-0-2)

\subsection*{Theory}
\noindent\textbf{Unit 1} \\
\textit{Ethical Hacking Fundamentals and Reconnaissance:} Cybersecurity kill chain and MITRE ATT\&CK framework, Legal and ethical considerations (ROE, NDA, scope agreements), Passive reconnaissance (OSINT, WHOIS, DNS enumeration, Shodan), Active reconnaissance (port scanning, service versioning, banner grabbing), Nmap scripting engine (NSE), Vulnerability scanning (Nessus, OpenVAS), Footprinting techniques and information leakage prevention.

\vspace{0.3cm}

\noindent\textbf{Unit 2} \\
\textit{System Hacking and Privilege Escalation:} Password cracking methodologies (dictionary, brute-force, rainbow tables), Hashcat and John the Ripper usage, Buffer overflow exploitation (stack/heap overflows), Privilege escalation techniques (Linux kernel exploits, Windows UAC bypass, SUID abuse), Password spraying and credential stuffing, Mimikatz for credential dumping, Lateral movement (Pass-the-Hash, Pass-the-Ticket).

\vspace{0.3cm}

\noindent\textbf{Unit 3} \\
\textit{Web Application Penetration Testing:} OWASP Top 10 vulnerabilities, SQL injection (blind, time-based, error-based), Cross-Site Scripting (XSS - reflected, stored, DOM-based), Cross-Site Request Forgery (CSRF/XSRF), Server-Side Request Forgery (SSRF), Authentication bypass techniques, Session management flaws, Business logic vulnerabilities, Burp Suite Professional workflow (proxy, repeater, intruder, scanner).

\vspace{0.3cm}

\noindent\textbf{Unit 4} \\
\textit{Network Penetration Testing and Wireless:} Network service exploitation (SMB, RDP, SSH weak credentials), Metasploit Framework (MSF) modules and payloads, Exploit development lifecycle, Post-exploitation modules (meterpreter, persistence, pivoting), Wireless attacks (WEP/WPA2 cracking, evil twin AP, KRACK), Man-in-the-Middle attacks (ARP poisoning, SSL stripping), Network device exploitation (router/switch configuration errors).

\vspace{0.3cm}

\noindent\textbf{Unit 5} \\
\textit{Advanced Pentesting and Reporting:} Active Directory attacks (Kerberoasting, AS-REP roasting, Golden/Silver tickets), Container and Kubernetes security testing, Cloud platform pentesting (IAM misconfigurations, S3 bucket enumeration), Red teaming methodologies and OPSEC, Penetration test reporting (executive summary, technical findings, risk ratings, remediation), Continuous security testing (DAST, SAST, IAST), Defensive strategies and blue teaming.

\newpage

\subsection*{Practical}
\noindent\textbf{Lab 1: Reconnaissance and Scanning} \\
Focuses on passive/active reconnaissance, Nmap scripting, vulnerability scanning with OpenVAS/Nessus.

\vspace{0.2cm}

\noindent\textbf{Lab 2: Password Cracking and Hash Analysis} \\
Focuses on cracking various hash types with Hashcat/John, creating custom wordlists and rules.

\vspace{0.2cm}

\noindent\textbf{Lab 3: Buffer Overflow Exploitation} \\
Focuses on stack/heap overflow exploitation, shellcode development, and ASLR bypass techniques.

\vspace{0.2cm}

\noindent\textbf{Lab 4: Metasploit Framework Mastery} \\
Focuses on MSF modules, payload customization, post-exploitation, and pivoting scenarios.

\vspace{0.2cm}

\noindent\textbf{Lab 5: SQL Injection Exploitation} \\
Focuses on discovering and exploiting SQLi vulnerabilities, database enumeration, privilege escalation.

\vspace{0.2cm}

\noindent\textbf{Lab 6: XSS and CSRF Attacks} \\
Focuses on reflected/stored XSS payloads, CSRF token bypass, DOM-based XSS exploitation.

\vspace{0.2cm}

\noindent\textbf{Lab 7: Burp Suite Web Pentesting} \\
Focuses on proxy interception, repeater manipulation, intruder fuzzing, and Burp extensions.

\vspace{0.2cm}

\noindent\textbf{Lab 8: Wireless and Network Attacks} \\
Focuses on WPA2 cracking, ARP poisoning, SSL stripping, and network service exploitation.

\vspace{0.2cm}

\noindent\textbf{Lab 9: Active Directory Domain Pentesting} \\
Focuses on Kerberoasting, DCSync attacks, and lateral movement in Windows environments.

\vspace{0.2cm}

\noindent\textbf{Lab 10: Complete Pentest Report} \\
Focuses on red team engagement simulation with executive/technical reporting and remediation roadmap.

\vspace{0.2cm}

\newpage
\rhead{Curriculum Schema}
